\begin{Exercise}[label={ex:picchi}]
    Si consideri una lista di interi. 
    Definiamo picchi della lista quei valori che sono preceduti e seguiti da valori
    più piccoli della loro metà. Il primo e l'ultimo elemento della lista non possono essere picchi.
    Scrivere una funzione che data una lista restituisce il numero di picchi della lista.
\end{Exercise}

\begin{Answer}[ref={ex:picchi}]
  	\lstinputlisting[style=cstylesmall]{src/picchi.c}
\end{Answer}